\chapter{Using the LeanFM Chat UI}
\label{ch:chat-ui-guide}

LeanFM's Chat UI takes a project from a problem statement to a maintained book,
checked requirements, a separately reviewed implementation mapping, source
files, and a running application. Work in the selected project's Chat pane;
use the other panes to inspect what the conversation actually produced. A
reply, a diagram, and a successful PDF build each answer different questions.
The two specification acceptance controls record the decisions that authorize
implementation.

This walkthrough uses the production interface with a disposable GoServer
project and a local test provider. The project contains real checked
specifications and runnable code. The displayed model name and token readings
are fixture data, not provider recommendations, prices, or performance results.
The screenshots contain no live account credentials or private conversations.
Their keyboard traces, accessibility snapshots, input revisions, and capture
boundaries are retained beside the images in \path{book/assets/ui-tour/}.

\newcommand{\UITourFigure}[4]{%
  \begin{figure}[htbp]
  \centering
  \AccessibleFigure{#3}{\includegraphics[width=\textwidth]{assets/ui-tour/#1.png}}%
  \caption{#2}\label{fig:ui-#1}
  \par\smallskip\footnotesize #4
  \end{figure}%
}

\section{Open the service and unlock your account}

Start LeanFM with \texttt{make serve} from its repository root, then open
\url{http://127.0.0.1:8080/}. Use the served LeanFM address when your browser is
on another machine. The application being developed has a separate runtime;
opening the LeanFM page does not itself start that application.

Select \textbf{Sign in} in the account controls. For an existing account, enter
your LeanFM username and password. This is the password that unlocks your
locally retained, encrypted provider keys. It is not your OpenAI or Anthropic
account password. After a LeanFM restart, sign in again to unlock the saved
keys; empty optional key fields do not mean those keys have disappeared.

\UITourFigure{01-sign-in}
  {Sign in unlocks the keys already retained for the LeanFM account.}
  {The actual LeanFM sign-in dialog shows username and password fields. Optional provider-key inputs and a password-change section precede the Log in button.}
  {Capture boundary: before LoginSubmitted. Both credential fields are empty.}

If you have no LeanFM account, choose \textbf{Create account}. Supply a username,
a password of at least eight characters, and at least one provider API key.
The form accepts OpenAI and Anthropic keys together. Put keys in these password
fields, never in a chat message, screenshot, source file, or project book.
Registration verifies the supplied provider keys before retaining them.

Once signed in, \textbf{API keys} opens the account's key editor. Re-enter the
LeanFM password to add or replace a key. Blank fields preserve retained keys.
Adding a second provider makes its models available without requiring a new
project. Replacing a key can invalidate older unlocked sessions. The optional
password-change fields require the current password and matching new values;
the service re-encrypts the retained keys and ends older sessions.

\UITourFigure{02-create-account}
  {The registration tab accepts one or both provider keys.}
  {The Create account tab is selected in the real authentication dialog. Username and password fields are empty. Both provider API key fields are empty.}
  {Capture boundary: before AccountCreateSubmitted. No registration or paid request is submitted for this image.}

\section{Choose the project before writing the prompt}

The left sidebar selects a project. Enter a descriptive name in \textbf{New
project name}, then choose \textbf{Add project} to create a separate directory
with its own identity, initial book, and specification paths.
Select the project and confirm its name in the workspace
header before sending a prompt. GoServer is an example project when present;
it is not a required dependency of creating another project.

The project owns its book, specifications, source, tests, generated views,
conversation, and evidence. Changing the selected tab changes the project
boundary. An active agent turn can continue while you inspect another project;
return to its tab to see its result. Delete a project only after checking the
name and the confirmation: deletion removes that project's complete directory
and stops its owned application.

The horizontal navigation groups the work by purpose:
\begin{description}[style=nextline]
\item[Chat] Describe the problem, answer consequential questions, and request
  concrete changes or checks.
\item[Book] Read or edit the durable design argument and preview its PDFs.
\item[Grammar and Tasks] Inspect the current project's generated behavioral
  views rather than relying on an illustration invented in a reply.
\item[Project memory] Review canonical artifact files and both acceptance gates.
\item[Source] Inspect the implementation files and tests actually written.
\end{description}
The sidebar's text-size controls change the interface scale. Project and phase
navigation preserve the mounted workspace state, including drafts and scroll
positions. Use \texttt{Tab} and \texttt{Shift+Tab} to move through controls and
\texttt{Enter} or \texttt{Space} to activate a focused button. A focused input
accepts ordinary typing; the multiline Chat composer has a separate Send button.

\UITourFigure{03-chat-ready}
  {Start in the selected project's Chat pane.}
  {The production workspace shows GoServer in the sidebar. Phase navigation and account identity appear above the app controls. The empty conversation has image attachment controls and a multiline composer with a model selector.}
  {Capture boundary: project opened, before ChatPromptSubmitted. The model belongs to the local documentation provider.}

\section{Send a useful prompt and understand the queue}

Begin with the problem and an observable outcome. For example:
\begin{quote}
People should be able to create named chat rooms, open a room, and post named
game challenges inside that room. Keep each room's conversation separate from
each challenge's conversation.
\end{quote}
Add a concrete failure case when it matters: another room must not receive the
message; a retry must not create a second challenge; an invalid action must show
an error. LeanFM uses the project files to resolve ordinary implementation
details. A consequential ambiguity or a missing acceptance decision can still
require your answer.

Select the model for the new prompt in \textbf{Model}. The header identifies its
provider and a fingerprint of the credential that will be used. The selection
is captured when a prompt enters the queue; changing the selector later does
not change a prompt already queued. LeanFM does not silently replace an
unavailable provider or key with another one.

Write in the composer and activate \textbf{Send}. You may keep typing and send
another prompt while the first is active. The visible queue runs in order, one
prompt at a time, so a later prompt receives the results of earlier work. Use a
queued item's \textbf{Remove} button to discard it before execution. Sending an
extra prompt does not interrupt the active one.

\textbf{Attach image} accepts PNG, JPEG, GIF, and static SVG, up to 2 MiB per
original and four selected attachments. A preview appears before submission;
remove an attachment there if it is unintended. GIF input is represented to
the model by its first frame. You can also submit an image without text.
\textbf{Copy reference} on a displayed image retains its stable reference for
a later prompt; copying the reference does not copy a credential or image bytes.

\UITourFigure{04-chat-running}
  {The active turn remains visible while the agent works.}
  {A real submitted fixture prompt appears in Chat. The Agent running banner names the captured model and shows elapsed time and the current stage. Cancel and the composer remain available.}
  {Capture boundary: ChatPromptSubmitted and ModelRequestIssued, before ModelResponseReturned. The local test provider deliberately holds the response for this capture.}

The running banner shows the selected model, elapsed time, and latest known
stage. Expand \textbf{Operational trace} to inspect real tool destinations,
checks, and failures. You may switch to Book, Tasks, Project memory, or Source
while the turn continues. The live trace is more useful than inferring progress
from the clock alone.

\textbf{Cancel} requests cancellation at a safe boundary. A provider request or
tool already in progress can finish; cancellation prevents subsequent work
after the server observes it. Wait for the active request to return before
assuming the task has stopped. A saved multi-batch task can expose
\textbf{Check progress}, \textbf{Resume task}, or \textbf{Restart from current
files}; these controls distinguish continuation from a fresh attempt.

\section{Read the result before retrying}

A completed assistant message retains its time, model, response duration, and
reported usage when applicable. The header's remaining-token figure is a local
session safety allowance. It is not the provider's monetary credit balance.
A native application launch uses the runtime controls and reports no provider
token charge of its own.

\UITourFigure{05-chat-response}
  {Review the answer and its evidence before taking the next action.}
  {The actual Chat pane shows the completed local-provider explanation of CreateRoom with its timestamp and usage. User-message Resend and message Delete controls appear beside the messages. The composer and model selector remain visible.}
  {Capture boundary: after ChatResponseRendered. The fixture reply explains existing checked behavior; it does not accept a gate or generate implementation files.}

\textbf{Resend} queues the selected user prompt again, retaining its attachment
references. It can cause a new provider request. Use it after understanding the
failure or changing the relevant input; it is not a promise that an uncertain
earlier request had no effect. \textbf{Delete} removes that conversation message.
It does not undo a committed source change, remove the book, or revoke an
accepted specification.

When a structured question appears, choose or type an answer in its form and
submit it. The answer closes the identified question and becomes a later queued
turn. Formal acceptance is handled by the native revision-bound controls in
Chat or Project memory. A successful compiler check, a model saying ``done,''
or a normal follow-up message does not accept a missing gate.

An error includes a concrete diagnostic and, where available, an expanded
\textbf{Diagnostic for Codex} with \textbf{Copy diagnostic}. Preserve that
diagnostic when reporting a problem. For a provider low-credit rejection, fix
that provider's billing condition or explicitly choose an available model and
credential before resending. LeanFM does not make an automatic paid retry for
that rejection. A local token-limit error, rate limit, incomplete response,
compiler failure, and network failure need different remedies; the operational
trace identifies which boundary failed.

\section{Keep the book and inspect the generated behavior}

Open \textbf{Book} to see the editable source beside its rendered PDF.
\textbf{Save source} saves the text; \textbf{Build \& preview} saves and builds
the reading copies. The status line distinguishes a saved source from a built
PDF. \textbf{Light PDF} and \textbf{Dark PDF} select equivalent renderings of
the same source. \textbf{Maximize PDF} expands the reading pane, and
\textbf{Restore split} returns to source and preview. The divider resizes the
two panes.

\UITourFigure{07-book}
  {The Book pane keeps editable source and the rendered reading copy together.}
  {The actual GoServer book editor shows saved LaTeX source. Save source and Build and preview controls appear above the editor. Light and Dark PDF selectors appear above the PDF preview pane.}
  {Capture boundary: saved source and an existing built PDF opened. This capture does not edit the book or accept a changed revision.}

The book records decisions and their explanation. \textbf{Tasks} opens the
views derived from the current typed specification. Select a task in the left
task list, then choose \textbf{Productions}, \textbf{States},
\textbf{Properties}, or \textbf{where\_}. Productions shows its request and
result alternatives; States shows its task state machine. Properties and
where\_ display the temporal requirements and contextual constraints. A blank
task set or missing grammar is a specification gap, not evidence that the
behavior exists.

\UITourFigure{08-task-states}
  {The room-creation state machine comes from the checked project specification.}
  {The Tasks pane selects Create a named chat room and the States tab. The generated diagram shows the request followed by server checks. Alternatives lead to RoomCreated success or an explicit RequestRejected result.}
  {Capture boundary: the current create\_room task rendered from its checked typed value. The companion requirement and task-catalog snapshots are retained with the capture.}

\textbf{Grammar} provides the nested production structure. Expand the relevant
production when you need to inspect composition and alternatives. Use Tasks
for a focused behavioral view and Grammar for its structure. Neither pane
should substitute another project's grammar or a separately drawn persuasive
picture for the selected project's checked model.

\section{Review both gates, then inspect the implementation}

Open \textbf{Project memory} to read the book, requirements, protobuf schema,
implementation mapping, and implementation-source summary. The status of a file
indicates whether it is absent, still bootstrap material, or current. A file's
presence alone does not establish that its content has passed the current
canonical check.

\UITourFigure{06-project-memory}
  {Project memory makes the durable files and acceptance state inspectable.}
  {The Project memory pane shows canonical artifact tabs and source inventory. Requirements and Implementation specification have distinct review cards. The demonstration project's actual current revisions are already accepted.}
  {Capture boundary: both recorded specification acceptances loaded. The screenshot does not press either acceptance button.}

The first gate accepts observable requirements and their derived outputs.
Review the task messages, outcomes, error behavior, boundaries, and generated
views; then use \textbf{Accept requirements} when the checked revision expresses
the intended product. The second gate accepts the separate implementation
specification and its traceability. Review how the chosen runtime, files,
transports, resource bounds, security controls, and tests realize those
requirements; then use \textbf{Accept implementation specification}.

The controls enable only for an eligible checked revision. Changing the
requirements or mapping can make an earlier acceptance stale. Return to the
appropriate review rather than treating the old approval as authorization for
different content. Once both current gates are accepted, ask Chat to implement
the actual source and run its checks. The accepted artifacts and existing
source remain the inputs; deleted conversation history must not become a
missing design dependency.

\UITourFigure{09-source}
  {Source displays the implementation files actually present in the project.}
  {The real Source pane lists GoServer source and test files and displays the selected file. This is an inspectable source inventory.}
  {Capture boundary: the accepted implementation's actual source tree opened.}

Use \textbf{Source} to inspect the changed files and tests. Request edits and
verification through Chat when needed; the Source pane is an inspection view.
Read the reported checks and their limitations. A successful startup confirms
that the app can start and answer its probes; it does not prove every game rule
or security property. GoServer's present slice provides rooms, named challenges,
and separate conversations. Playable boards and the larger Go rules remain
deferred.

\section{Start, open, and stop the app}

Use \textbf{Start app} in the selected project's application bar after its
current specifications are accepted and its source is runnable. LeanFM uses
the fixed application contract, waits for actual readiness and a page response,
and keeps the owned process running. The bar changes to \textbf{running},
\textbf{Open app} becomes a link, and Chat receives the same verified launch
link. These native controls do not require a model request to perform the
launch.

\UITourFigure{10-app-ready}
  {A verified running instance produces Open app in both the bar and Chat.}
  {The actual project workspace shows App running with Start app and Stop app controls and a served Open app link. A native runtime chat message also provides a clickable Open app link.}
  {Capture boundary: ApplicationReady after real startup and page probes. The link names the served preview, not a private loopback endpoint.}

\textbf{Open app} opens the served preview. The app runs inside its own
filesystem, process, network, and browser-authority boundary. For a remote
LeanFM browser, use this link rather than guessing a server-local port. Apps
need relative or explicitly prefix-compatible asset and API URLs for the
preview; the runtime reports an unsupported contract instead of publishing a
link that cannot work.

\UITourFigure{11-app-preview}
  {The running GoServer preview supports real room and challenge operations.}
  {The actual isolated GoServer browser shows a created room and a named challenge. Separate room and challenge conversation areas include a message delivered through the preview API.}
  {Capture boundary: RoomCreated, GameCreated, and GameChat observed through the real preview route. The app remains a local in-memory inspection slice.}

If source changes, the runtime becomes \textbf{stale}. Use \textbf{Restart app}
to run the new revision. Repeated Start on a current running revision reuses its
instance. \textbf{Stop app} ends only the selected project's owned instance and
removes its ready link. Other projects and unrelated processes are left alone.
The separate terminating launch check is validation evidence; its old receipt
is not a persistent Open app link.

All HTTP applications expose the same targets from \texttt{source/}:
\begin{center}
\begin{tabular}{ll}
\toprule
Target & Purpose \\
\midrule
\texttt{make build} & Build the application. \\
\texttt{make run} & Run in the foreground for supervision. \\
\texttt{make serve} & Alias for \texttt{run}. \\
\texttt{make stop} & Stop only this project's owned app. \\
\texttt{make test} & Execute the application's tests. \\
\texttt{make check} & Run its proportional verification. \\
\texttt{make clean} & Remove its disposable build outputs. \\
\bottomrule
\end{tabular}
\end{center}
Missing targets, invalid launch profiles, capacity limits, and failed readiness
produce explicit errors. Review the app status and diagnostic before retrying;
an occupied listener is not stopped to make room for the new app.

\section{Reproduce and review the screenshots}

Run \texttt{make ui-snapshots} from the LeanFM repository after building the
current server and UI. The runner uses the real production bundle, a disposable
project replica, an isolated network, and a deterministic local provider. It
drives named controls by keyboard and records the focused control, accessible
roles and states, visible message IDs, viewport, browser version, source
digests, and exact capture boundary. No live provider key or production
conversation is required. The GoServer example is needed to regenerate this
particular tour; the checked-in capture inputs and images allow the book to
build independently of that optional project directory.

Each image has a neighboring \texttt{.accessibility.json} and
\texttt{.keyboard.json} file. \path{book/assets/ui-tour/manifest.json} binds them
to the prompt and source revisions. The book's preflight check rejects missing,
altered, or stale required captures. Alternative text in both PDF editions
describes the control state, so a screenshot is never the only account of the
interaction.
